\chapter{Introduction} \label{ch:Einleitung}

\section{Task and Motivation}
\label{sec:AuM}
In this experiment, the generation and analysis of X-ray radiation are investigated. The aim is, on the one hand, to determine fundamental constants of physics experimentally and, on the other hand, to analyze the structure of crystals using diffraction methods. Specifically, Planck's quantum of action $h$ is to be determined from the short-wavelength end of the bremsstrahlung spectrum as well as from the threshold voltage at characteristic wavelength. Furthermore, the wavelengths of the characteristic $K_\alpha$ and $K_\beta$ lines of a molybdenum anode are to be determined. By using lithium fluoride (LiF) and sodium chloride (NaCl) as diffraction crystals, the connection between crystal lattice structure and macroscopic quantities such as the Avogadro number $N_A$ can also be established.

\section{Physical Basis}
\label{sec:PhyBasics}

\subsection{Generation of X-ray Radiation}
The generation of X-ray radiation in this experiment takes place via an X-ray tube with a molybdenum anode and a heated filament cathode. Inside an evacuated glass vessel, electrons are accelerated by an electric field between the cathode and the anode. The kinetic energy $E$ of the electrons results from the applied accelerating voltage $U$:

\eq{energie}{
    E = e U
    }
where $e$ represents the elementary charge. 
Upon impact with the anode, the electrons are decelerated in the Coulomb field of the atomic nuclei. This braking process leads to the emission of a continuous spectrum of electromagnetic radiation, known as the bremsstrahlung spectrum.
\img[0.7]{X_Ray_Sensor}{Creating X-Ray. \cite{skriptPAP22}}

The maximum energy of an emitted photon is reached when an electron converts its entire kinetic energy into a single photon. This results in a minimum possible wavelength $\lambda_{gr}$, the cutoff wavelength (or short-wavelength limit), which is described by conservation of energy:
\eq{grenzwelle}{
    e U = h \nu_{gr} = \frac{hc}{\lambda_{gr}}
    }
Here, $h$ is Planck's quantum of action and $c$ is the speed of light. In addition to the continuous spectrum, a characteristic line spectrum is generated if the voltage is sufficiently high. This occurs when electrons are knocked out of inner shells of the anode material and the resulting vacancies are filled by electrons from higher shells. For the transition from the $n$-th to the $m$-th shell, the energy $E_{n \to m}$ can be estimated according to Moseley's Law:
\eq{mosley}{
    E_{n \to m} = hc R_\infty (Z - A)^2 \left( \frac{1}{m^2} - \frac{1}{n^2} \right)
    }
Here, $R_\infty$ denotes the Rydberg constant, $Z$ the atomic number of the anode material (here molybdenum), and $A$ a shielding constant, which is approximately $A \approx 1$ for $K_\alpha$ radiation ($n=2 \to m=1$).

\subsection{Bragg Reflection}
\wrap[0.5]{R}{B_scattering}{Idea of Bragg scattering. \cite{skriptPAP22}}
To measure the wavelengths of the generated X-ray radiation, crystals are used as diffraction gratings. Since the wavelengths of X-ray radiation are of the same order of magnitude as the interatomic distances in crystals, constructive interference occurs at the lattice planes of the crystal. This phenomenon is called Bragg reflection. An incident ray is reflected at various parallel lattice planes. For these reflected rays to overlap constructively, the path difference $\Delta s$ must be an integer multiple of the wavelength $\lambda$. Bragg's Law describes this relationship:
\eq{bragg}{
    2 d \sin(\vartheta) = n \lambda, \quad n \in \mathbb{N}
}
Here, $d$ is the distance between adjacent lattice planes, $\vartheta$ is the angle of incidence (glancing angle), and $n$ is the order of reflection. By varying the angle $\vartheta$ and measuring the intensity, the X-ray spectrum can be recorded.


\subsection{Crystals and Lattice Structure}
The LiF and NaCl crystals used in the experiment both possess a face-centered cubic structure. In this structure, a unit cell repeats periodically in space. There are four atoms (or ion pairs) per unit cell. For a cut along the principal planes, the lattice plane spacing $d$ in terms of the lattice constant $a$ is given by:
\eq{netzebene}
{
    d = \frac{a}{2}
    }
If the wavelength $\lambda$ is known, the lattice constant $a$ can be determined from the measured angle $\vartheta$. Conversely, a known lattice constant enables the determination of microscopic quantities from macroscopically measurable properties. The Avogadro number $N_A$, which indicates the number of particles per mole, can be calculated via the volume of the unit cell. The volume of a unit cell is $V = (2d)^3$. With the molar mass $M_{mol}$ and the density $\rho$ of the crystal, it follows that:
\eq{avogadro}{
    N_A = \frac{4 V_{mol}}{(2d)^3} = \frac{4 M_{mol}}{\rho (2d)^3} = \frac{1}{2} \frac{M_{mol}}{\rho d^3}
    }
This relationship allows the Avogadro number to be determined indirectly via X-ray diffraction.
In the evaluation, the statistical nature of radioactive or count-rate-like events is also taken into account. Since the count rate $M$ follows a Poisson process, the standard error of a single measurement is:
\eq{fehler_zahl}{
    \sigma = \sqrt{M}
    }
This value is used to incorporate the uncertainty of the measured intensities into the error propagation.
\img{Lattice_Structure}{Structure of the crystals in use; Here NaCl. \cite{skriptPAP22}}